\begin{example} 
    \label{ex:ax-b}
    Sea $(x_n)_{n\geq1}$ una sucesión equidistribuida módulo $1$, sea $a\in \mathbb{Z}\setminus\{0\}$ y sea $b\in\mathbb{R}$.
    Entonces la sucesión $(ax_n+b)_{n\geq1}$ es equidistribuida módulo $1$.
\end{example}

\begin{resolve}
    Definimos $y_n = ax_n+b$. Para cada $k \in \mathbb{Z}\setminus\{0\}$ se tiene
    \[
        \left|\frac{1}{N}\sum_{n=1}^N e^{2\pi i k y_n} \right|
        = \left|\frac{1}{N}\sum_{n=1}^N e^{2\pi i k (ax_n + b)} \right|
        = \left|\frac{1}{N}e^{2\pi i k b}\sum_{n=1}^N e^{2\pi i (ka) x_n} \right|
        =\left|\frac{1}{N}\sum_{n=1}^N e^{2\pi i (ka) x_n} \right|.
    \]
    
    Dado que $ka \in \mathbb{Z}\setminus\{0\}$ y $(x_n)$ es equidistribuida, el criterio de Weyl implica que
    \[
        \lim_{N\to\infty} \frac{1}{N}\sum_{n=1}^N e^{2\pi i (ka) x_n} = 0.
    \]
    
    Por consiguiente,
    \[
        \lim_{N\to\infty} \frac{1}{N}\sum_{n=1}^N e^{2\pi i k y_n}  = 0.
    \]
    y por el criterio de Weyl, la sucesión $(y_n)$ es equidistribuida módulo $1$.
\end{resolve}

Como consecuencia, la proposición \ref{prop:secuencia_mas_const} es un corolario de este resultado.